\chapter{Conclusions and Future Work} \label{chap:ch5}

\section{Summary of the Developed Work}

This dissertation presented the design, physical implementation, and verification of a 10 Gbps \ac{CML} transmitter targeting the \ac{USB}~3.2 Gen~2 electrical requirements. The work followed the transmitter from specification analysis to transistor-level design, layout, parasitic extraction, and post-layout verification in the GlobalFoundries \SI{55}{\nano\meter} LPx \ac{CMOS} technology.

The design was motivated by the \ac{SI} limitations that arise in multi-gigabit serial links. At 10 Gbps, the channel introduces substantial frequency-dependent attenuation and therefore significant \ac{ISI}. The background study in Chapter~\ref{chap:ch2} established the relationship between channel loss, eye closure, jitter, \ac{TX}-line behaviour, and transmitter-side equalization, and provided the basis for selecting a current-mode, three-tap \ac{FFE} architecture.

The proposed transmitter, described in Chapter~\ref{chap:ch3}, implements independently configurable pre-cursor, main-cursor, and post-cursor taps. Their currents are generated by segmented current-mirror banks and are digitally programmable, allowing the transmitted waveform to be adjusted without changing the underlying analog topology. The required tap coefficients were derived directly from the \ac{USB} preshoot and de-emphasis definitions and converted into current targets and enable-code mappings. A programmable termination network was integrated in the same modular structure so that output impedance can also be adjusted across process and resistor variations.

A major part of the work concerned the physical implementation of the transmitter. The tap unit was developed as a modular, abuttable layout cell using common-centroid placement, symmetric routing, dummy devices, continuous guard-ring structures, and routing strategies intended to reduce mismatch and layout-dependent effects such as \ac{WPE} and \ac{LOD}. The complete extracted transmitter, including the tap arrays, biasing circuitry, and termination network, occupies \SI{176}{\micro\meter} $\times$ \SI{82}{\micro\meter}.

Chapter~\ref{chap:ch4} evaluated the transmitter bfeore and after parasitic extraction and across the relevant \ac{PVT} and resistor corners. The results show that the physical implementation has a significant effect on the electrical behaviour of the driver. For the representative EQ3 configuration at the \ac{TT}, Res\_nom corner, the post-extraction preshoot and de-emphasis magnitudes are 3.121~dB and 6.114~dB, respectively. The output rise and fall times increase to 42.1~ps and 35.6~ps, compared with 27.3~ps and 22.8~ps before extraction.

Despite the slower output transitions, the received eye remains widely open after propagation through the \ac{USB-IF} reference channel. EQ3 produces an \ac{RX} eye height of \SI{322}{\milli\volt} and a width of \SI{84.7}{\pico\second}, while EQ4 produces \SI{323}{\milli\volt} and \SI{84.3}{\pico\second}, respectively. The two presets therefore provide essentially equivalent post-extraction performance. EQ4 gives a marginally larger vertical opening, whereas EQ3 provides both a marginally larger horizontal opening and a slightly smaller crossover width, \SI{15.6}{\pico\second} compared with \SI{16.0}{\pico\second}. EQ3 was consequently retained as the representative configuration for the detailed corner analysis. The nominal post-extraction power consumption of EQ3 is \SI{24.42}{\milli\watt} at \ac{TT} and Res\_nom.

The pre- and post-extraction comparison demonstrates that layout parasitics cannot be treated as a secondary effect at this data rate. At the representative \ac{TT} operating point, the rise and fall times increase by approximately 54.2\% and 56.1\%, respectively, showing a comparable degradation of both output transitions. The EQ3 de-emphasis magnitude decreases by approximately 7.8\%, and the termination-bank setting required to obtain the nominal output impedance also changes after extraction. These results demonstrate that parasitic extraction affects not only the transition speed but also the effective equalization weighting and termination operating point.

The equalization results further demonstrate the importance of the segmented architecture. The nominal analytical tap mappings are not preserved uniformly after extraction, with de-emphasis becoming particularly sensitive at the \ac{SS} corner. The EQ1 retuning study shows that the final compliant mapping is itself \ac{PVT}-dependent: \ac{TT} and \ac{FF} are corrected using the 7-102-16 pre-cursor/main-cursor/post-cursor enable mapping, whereas \ac{SS} requires the substantially different 4-82-20 mapping. The programmable current banks therefore provide sufficient tuning range to recover valid equalization operating points, but a single universal post-layout equalization code is not sufficient across the complete evaluated \ac{PVT} space.

When compared with representative \ac{CML} transmitters reported in the literature, the developed architecture occupies a different point in the power--equalization--robustness trade-off. Earlier current-mode transmitters demonstrate experimentally the effectiveness of multi-tap \ac{FFE}, while hybrid and capacitive equalization techniques show that lower transmitter power can be obtained by reducing static equalization current or by operating with very small channel swings. The contribution of the present work is instead the combination of a three-tap current-domain equalizer with programmable termination and sufficient digital adjustment range to compensate the changes introduced by physical implementation and \ac{PVT} variation. The state-of-the-art comparison in Section~\ref{sec:sota_comparison} therefore supports the positioning of the proposed transmitter as a configurable, layout-aware 10 Gbps \ac{CML} implementation rather than as an absolute minimum-power or maximum-data-rate design.

Overall, the work demonstrates that the proposed segmented \ac{CML} architecture is capable of operating at the target data rate and provides sufficient digital control range to obtain \ac{USB}~3.2 Gen~2-compatible output swing, equalization, output impedance, and eye-opening operating points after physical implementation. The results validate the usefulness of both the programmable circuit architecture and the modular layout methodology, while also showing that calibration becomes an important part of achieving robust operation across \ac{PVT}. Complete compliance validation still requires the normative jitter and \ac{BER} methodology together with experimental characterization in silicon.

\section{Fulfillment of the Objectives}
The objectives introduced in Chapter~\ref{chap:ch1} were addressed throughout the dissertation as summarized in Table~\ref{tab:objectives_summary}. In particular, the work progressed from the theoretical analysis of the link and equalization problem to a complete post-layout transmitter implementation, allowing the original specifications to be evaluated against extracted circuit behaviour rather than only against an ideal schematic.

\begin{table}[!ht]
\centering
\caption{Summary of the fulfillment of the dissertation objectives.}
\label{tab:objectives_summary}
\begin{tabular}{p{0.56\textwidth} p{0.2\textwidth} p{0.20\textwidth}}
\hline
\textbf{Objective} & \textbf{Status} & \textbf{Reference} \\
\hline
Study the \ac{SI} mechanisms affecting high-speed serial links, including \ac{ISI}, jitter, channel loss, and equalization. &
Completed &
Chapter~\ref{chap:ch2} \\
Develop a segmented \ac{CML} transmitter implementing a three-tap \ac{FFE}. &
Completed &
Chapter~\ref{chap:ch3} \\
Derive tap coefficients and current requirements from the \ac{USB}~3.2 Gen~2 \ac{PS} and \ac{DE} targets. &
Completed &
Section~\ref{sec:tx_equalization} \\
Design the differential pairs, current mirrors, biasing, and programmable current sources required by the transmitter. &
Completed &
Chapter~\ref{chap:ch3} \\
Develop a modular and abuttable physical implementation accounting for layout-dependent effects and differential routing symmetry. &
Completed &
Section~\ref{sec:layout} \\
Verify the \ac{TX} before and after parasitic extraction across \ac{PVT} and resistor corners. &
Completed &
Chapter~\ref{chap:ch4} \\
Evaluate the implemented \ac{TX} against the relevant \ac{USB}~3.2 Gen~2 electrical requirements. &
Completed with digital trim &
Section~\ref{sec:usb32_compliance}\\
\hline
\end{tabular}
\end{table}

The final objective is qualified as ``completed with digital trim'' because the post-extraction results demonstrate specification-compatible operating points, but the required digital configuration is not universal across all evaluated conditions. This is particularly evident in the equalization and termination results.

For EQ1, after extraction, the same current-mirror allocation does not provide adequate equalization at every PVT corner. The final selected mapping is 7-102-16 for both \ac{TT} and \ac{FF}, whereas the \ac{SS} corner requires 4-82-20. Similarly, although the termination network does not reproduce the complete 72~$\Omega$ to 120~$\Omega$ tuning interval for every combination of transistor and resistor corners, at least one compliant termination setting is available for each evaluated condition.

The objective is therefore considered fulfilled at the level of programmable post-layout verification: the implemented architecture provides sufficient adjustment range to recover the evaluated electrical requirements. A practical silicon implementation would, however, require the appropriate equalization and termination codes to be determined through calibration or pre-characterized operating tables.

\section{Limitations}

Although the post-extraction results demonstrate the functionality of the proposed transmitter and the ability of its digital programmability to recover suitable operating points, several limitations remain in the present implementation and verification methodology.

The first limitation is that the results presented in this work are entirely simulation-based. The transmitter was verified at schematic and post-extraction levels, but it was not fabricated and experimentally characterized. Consequently, the reported performance does not include measurement uncertainty, fabrication mismatch, package effects, or other silicon-specific non-idealities that can only be quantified after tape-out. The presented results should therefore be interpreted as post-layout design verification rather than as measured compliance of a fabricated transmitter.

A second limitation is the absence of a complete jitter-compliance characterization. The eye crossover width was used as an indicator of the received transition dispersion, and the strongest equalization presets show a substantial reduction in crossover width after extraction. At the \ac{TT} corner, the EQ3 crossover width decreases from \SI{45.1}{\pico\second} before extraction to \SI{15.6}{\pico\second} after extraction, while the EQ4 crossover width decreases from \SI{49.7}{\pico\second} to \SI{16.0}{\pico\second}. The post-extraction EQ3 crossover width also remains relatively stable across the evaluated \ac{PVT} corners, with values of \SI{15.6}{\pico\second}, \SI{16.9}{\pico\second}, and \SI{15.1}{\pico\second} for \ac{TT}, \ac{FF}, and \ac{SS}, respectively. These measurements indicate a comparatively concentrated distribution of the received transition crossings after extraction. 

However, eye crossover width is only an indicator of timing dispersion and cannot be directly equated with the normative jitter quantities. The complete normative \ac{TX}-eye and jitter procedure, including the required compliance patterns, reference receiver equalization, \ac{DJ}/\ac{RJ}/\ac{TJ} characterization, measurement over $10^6$ consecutive \ac{UI}, and jitter extrapolation to a \ac{BER} of $10^{-12}$, was not performed.

A third limitation is the dependence of post-extraction equalization accuracy on \ac{PVT} and layout parasitics. The initial tap-current allocations were derived from ideal relationships between the pre-cursor, main-cursor, and post-cursor coefficients, whereas the extracted implementation modifies their effective relative weighting. At the \ac{TT} and \ac{FF} corners, preshoot generally remains comparatively close to its target, while de-emphasis exhibits a larger reduction. This effect becomes substantially more pronounced at the SS corner, particularly for the stronger equalization presets. For example, the untuned EQ1 configuration produces 1.41~dB of de-emphasis magnitude at SS, compared with its nominal 2.1~dB target, corresponding to a reduction of approximately 32.9\%. Even larger relative \ac{DE} errors are observed for some of the stronger presets at this corner.

The segmented current-source architecture provides sufficient tuning range to correct this behaviour, but the EQ1 tuning results demonstrate that the required correction is itself corner-dependent. The selected post-extraction mapping is 7-102-16 for both \ac{TT} and \ac{FF}, whereas the \ac{SS} corner requires the substantially different 4-82-20 mapping. Consequently, a single fixed equalization code cannot be assumed to provide the required behaviour across the complete evaluated \ac{PVT} space. The present implementation does not include an automatic equalization-calibration loop or a stored corner-dependent configuration procedure, and therefore relies on the appropriate digital code being selected externally.

The available simulations also do not uniquely identify a single physical mechanism or individual \ac{FFE} tap as the cause of the post-extraction equalization error. Parasitic extraction simultaneously modifies the resistance and capacitance associated with the differential signal paths, the current-source routing, and the shared output nodes. The observed equalization changes should therefore be interpreted as the combined effect of the physical implementation on the relative tap contributions. A more detailed sensitivity analysis in which the parasitic contributions of the individual tap paths are introduced or removed independently would be required to quantify their separate contributions.

The programmable termination network presents a related limitation. Its segmented implementation provides sufficient flexibility to obtain an output impedance within the specified range for every evaluated combination of \ac{PVT} and resistor corner, but the digital code corresponding to a target impedance changes substantially after parasitic extraction and across resistor corners. Furthermore, the complete \SIrange{72}{120}{\ohm} tuning interval is not available at every combined corner, even though at least one compliant operating point is available in each evaluated case. A single fixed termination code therefore cannot guarantee the target impedance over all operating conditions. In a fabricated implementation, the appropriate code would need to be determined from the actual device characteristics or selected through a calibration procedure.

The physical implementation also introduces a trade-off between modularity, programmability, routing parasitics, and occupied area. The repeated tap-unit architecture simplifies layout reuse, matching, and digital adjustment, but requires a comparatively large number of devices, control signals, and parallel interconnects. The post-extraction increase in transition times illustrates the electrical cost of this physical complexity. At the representative \ac{TT} condition, the rise time increases from 27.3~ps to 42.1~ps, while the fall time increases from 22.8~ps to 35.6~ps, corresponding to increases of approximately 54.2\% and 56.1\%, respectively. Further reduction of routing resistance and capacitance could improve the high-speed response, but such changes must be balanced against area, routing congestion, \ac{EM} constraints, and preservation of the modular layout methodology.

Finally, the \ac{SI} verification was performed using one specific \ac{USB-IF} reference channel. The selected channel is intentionally demanding, with \SI{21.6}{\decibel} insertion loss at the \SI{5}{\giga\hertz} Nyquist frequency, and provides a useful benchmark for evaluating the equalization capability of the transmitter. Nevertheless, it represents only one link condition rather than the complete range of channels that can occur in a practical system. The equalization setting that maximizes the received eye can change with channel loss, package parasitics, connector discontinuities, and receiver characteristics. Verification over a broader set of channel models and, ultimately, with a complete transmitter--channel--receiver link would therefore be required to establish the robustness of the proposed equalization strategy under a wider range of system conditions.

\section{Future Work}

The most direct continuation of this work is the fabrication and experimental characterization of the proposed transmitter. A silicon prototype would allow the simulated current-source accuracy, output swing, differential output impedance, transition times, equalization settings, eye diagrams, and power consumption to be compared directly with measured results. This would also make it possible to quantify the correlation between the schematic, post-extraction, and fabricated implementations and to identify additional non-idealities associated with device mismatch, package parasitics, supply noise, and measurement interfaces that are not captured completely by the present simulation methodology.

A fabricated implementation should also undergo a complete physical-reliability sign-off and validation. In addition to functional and signal-integrity measurements, the final chip should be evaluated for \ac{EM} and supply-network IR drop under the worst expected operating conditions. This is particularly relevant for the transmitter output stage, programmable current-source banks, termination network, and their associated supply and ground interconnects, where relatively large currents are concentrated in a limited routing area. Excessive current density can accelerate electromigration and compromise long-term interconnect reliability, while resistive voltage drops in the power-distribution network can modify the effective bias conditions, output swing, tap-current ratios, and consequently the equalization behaviour of the transmitter.

Such verification should include current-density checks in metal segments and vias, static and dynamic IR-drop analysis, and thermal or self-heating assessment under representative and worst-case switching activity. These analyses should be performed across the relevant \ac{PVT} conditions and should account for the simultaneous activation of the current-source and termination segments associated with the most demanding transmitter configurations. Where necessary, the power-routing width, via multiplicity, and local supply distribution should be reinforced to satisfy the technology reliability limits. Combined with silicon measurements and, where applicable, accelerated stress testing, these checks would provide confidence that the transmitter not only meets its electrical specifications at initial operation but also maintains them throughout the intended operating lifetime.

Experimental characterization should also include a complete \ac{USB}~3.2 Gen~2 transmitter compliance procedure. In particular, the normative decomposition of the transmitter jitter into \ac{DJ}, \ac{RJ}, and \ac{TJ} should be performed using the appropriate reference receiver and statistical methodology defined by the standard. A direct \ac{BER}-based evaluation should likewise be included. Although the eye crossover measurements obtained in this work provide a useful indication of the received transition dispersion, they cannot replace these normative jitter and error-rate measurements. Silicon characterization would therefore provide the final step required to determine whether the post-layout electrical compatibility demonstrated in this dissertation translates into complete experimental compliance.

A particularly important extension is the development of a calibration strategy for both equalization and termination. The post-extraction EQ1 results demonstrate that the required current-source configuration is \ac{PVT}-dependent. The selected compliant mapping is 7-102-16 at both \ac{TT} and \ac{FF}, whereas the \ac{SS} corner requires the substantially different 4-82-20 mapping. The existing segmented architecture already provides the actuation mechanism required for calibration, since the pre-cursor, main-cursor, and post-cursor current-source mirrors can be controlled independently.

A practical implementation could therefore employ either a start-up calibration procedure or a set of pre-characterized configuration tables to select the appropriate tap-current codes according to the actual operating conditions. The same calibration process could determine the termination-bank code required to place the differential output impedance within its target range. Ultimately, an automatic calibration loop would remove the need to assume a single universal equalization or termination code and would make the transmitter substantially more robust to process, supply-voltage, temperature, and passive-device variation.

The post-extraction results also identify opportunities for further physical optimization. The substantial increase in output rise and fall times shows that the parasitic resistance and capacitance associated with the high-speed differential paths remain important limitations of the current layout. Furthermore, the change in equalization accuracy after extraction indicates that the physical implementation modifies the relative contribution of the three \ac{FFE} taps. Future layout iterations should therefore aim to reduce and balance the parasitic loading of the differential output and tap-current paths, while also reducing the series resistance associated with the programmable termination network.

These modifications should be guided by a more detailed parasitic sensitivity analysis. In the present work, the extracted network includes the parasitic contributions of all tap and output paths simultaneously, making it difficult to identify the individual contribution of each physical mechanism to the observed equalization error. A future analysis could selectively introduce the extracted resistance and capacitance associated with the pre-cursor, main-cursor, post-cursor, and shared output paths in order to quantify their separate impact on preshoot, de-emphasis, transition time, and received eye opening. This would allow subsequent layout changes to target the dominant parasitic contributions rather than relying on qualitative assumptions regarding the origin of the post-extraction error.

The current-source segmentation and digital activation strategy could also be investigated further. Alternative distributions of enabled current-mirror segments across the available tap units could be compared to determine whether a more spatially distributed activation sequence improves current matching or reduces sensitivity to local layout-dependent effects. Such an investigation should be supported by extracted and mismatch-aware simulations, since the present results are not sufficient to conclude whether activation ordering is a dominant source of the measured tap-current error. Any improvement would also need to be balanced against the additional decoder complexity, routing density, and layout area required by a more elaborate control scheme.

Another useful extension would be to incorporate post-extraction equalization optimization directly into the transmitter design flow. In the present methodology, the initial digital preset mappings are derived from analytical tap coefficients and subsequently corrected after extraction. The strong corner-dependence observed for EQ1 demonstrates that the final physical implementation can require substantially different configurations from those predicted by the ideal current relationships.

A future optimization flow could therefore include the extracted transmitter, the target channel, and the relevant \ac{PVT} condition directly in an automated search for the digital current-source and termination codes. The optimization criterion should not be limited to minimizing the numerical error relative to the nominal preshoot and de-emphasis targets. As demonstrated by the EQ1 tuning procedure, the configuration having the smallest target error is not necessarily the preferred configuration if one of the resulting quantities lies outside its normative range. The optimization should therefore treat the \ac{USB} preshoot, de-emphasis, swing, and impedance intervals as explicit constraints while optimizing a system-level metric such as received eye height, eye width, power consumption, or a combination of these quantities.

The channel space considered during verification could likewise be expanded. The present work uses one demanding \ac{USB-IF} reference channel with \SI{21.6}{\decibel} insertion loss at the \SI{5}{\giga\hertz} Nyquist frequency. Future verification should include channels with different insertion-loss profiles, impedance discontinuities, connector and package parasitics, and different overall link lengths. This would allow the required equalization settings to be characterized as a function of the actual channel rather than for a single reference case, and would provide a more complete evaluation of the useful tuning range offered by the segmented transmitter.

Finally, the transmitter could be extended toward a complete \ac{PHY} implementation. Integration with the preceding high-speed data path, digital control and calibration logic, realistic package and interface models, and a representative receiver would enable end-to-end link verification. Such a platform would allow the interaction between transmitter equalization, channel loss, receiver equalization, jitter, and \ac{BER} to be evaluated jointly, rather than treating the transmitter as an isolated block.

The next development stage is therefore centered on converting the post-layout-verified transmitter into a calibrated, reliability-verified, and experimentally validated system. The results of this dissertation indicate that the segmented architecture provides sufficient adjustment range after physical implementation, but they also show that this programmability must be used actively to compensate for layout- and \ac{PVT}-dependent behaviour. Automatic calibration, targeted physical optimization, \ac{EM} and IR-drop verification, broader channel characterization, and silicon measurement therefore constitute the principal steps toward a complete and robust \ac{USB}~3.2 Gen~2 transmitter implementation.
